\section{Methodology}
\label{sec:methodology}

Our framework, \textbf{Mem0-Cognitive}, reframes memory management as an \textit{Active Inference} process. Instead of passively storing all interactions, the system continuously predicts the future utility of each memory trace and adapts its retention strategy accordingly. Figure~\ref{fig:architecture} illustrates the overall architecture, which comprises three core modules: (1) Emotion-Weighted Forgetting, (2) Offline Sleep Consolidation, and (3) Meta-Cognitive Adaptation.

\subsection{Problem Formulation}
Let $\mathcal{M}_t = \{m_1, m_2, \dots, m_n\}$ denote the memory store at time $t$, where each memory $m_i$ is a tuple $(c_i, e_i, ts_i, \theta_i)$ consisting of content $c_i$, embedding $e_i$, timestamp $ts_i$, and cognitive parameters $\theta_i$. The goal is to learn a policy $\pi_\phi$ that optimizes memory operations (store, update, compress, delete) to maximize long-term retrieval utility:
\begin{equation}
    \max_\phi \mathbb{E}_{t \sim T} \left[ \sum_{i=1}^{|\mathcal{M}_t|} U(m_i, q_t) - \lambda \cdot \text{Cost}(\mathcal{M}_t) \right]
\end{equation}
where $U(m_i, q_t)$ is the utility of memory $m_i$ for query $q_t$, and $\text{Cost}(\cdot)$ penalizes excessive memory size (token overhead).

\subsection{Emotion-Weighted Forgetting with Salience Gate}
\label{subsec:forgetting}

Classic Ebbinghaus forgetting models describe memory retention as $R(t) = e^{-t/S}$, where $S$ is a constant strength parameter. However, this fails to account for emotional salience, which in human cognition significantly modulates forgetting rates. We introduce a \textbf{Salience Gate} mechanism $\mathcal{G}$ that dynamically adjusts the effective decay rate based on extracted emotional intensity.

Given a memory $m_i$ with base strength $S_{base}$ and emotional valence $E_i \in [0, 1]$ (extracted via zero-shot LLM prompting; see Appendix~\ref{app:prompts}), the effective retention strength is:
\begin{equation}
    S_{\text{eff}}(m_i) = S_{\text{base}} \cdot (1 + \lambda \cdot E_i)
\end{equation}
where $\lambda \in [0, 2]$ controls the \textit{emotional inertia} (higher $\lambda$ implies stronger resistance to forgetting for emotional events). The retention probability at time $\Delta t$ after encoding is then:
\begin{equation}
    \mathcal{R}_{\text{affective}}(m_i, \Delta t) = \exp\left(-\frac{\Delta t}{S_{\text{base}} \cdot (1 + \lambda \cdot E_i)}\right)
\end{equation}

To prevent catastrophic forgetting of emotionally charged but rarely accessed memories (e.g., user traumas or core identity statements), we introduce a \textbf{Salience Gate} threshold $\tau_{salience}$. Memories with $E_i > \tau_{salience}$ are exempt from automatic deletion and instead undergo compression. This ensures that high-salience information persists even with low access frequency.

The final \textit{Affective Retention Score} combines temporal decay with access frequency $f_i$ and recency weight $w_{recency}$:
\begin{equation}
    \psi_{\text{retention}}(m_i) = \alpha \cdot \mathcal{R}_{\text{affective}}(m_i, \Delta t) + \beta \cdot \text{norm}(f_i) + \gamma \cdot w_{recency}
\end{equation}
where $\alpha, \beta, \gamma$ are hyperparameters balancing the contributions ($\alpha=0.5, \beta=0.3, \gamma=0.2$ in our experiments).

\subsection{Offline Sleep Consolidation as Schema Induction}
\label{subsec:consolidation}

Inspired by hippocampal-neocortical dialogue during NREM sleep~\cite{mcclelland1995complementary}, our consolidation engine operates as an \textit{offline schema induction module}. Unlike online reflection mechanisms that block user interaction, sleep consolidation runs asynchronously during system idle periods.

Given a cluster of episodic memories $\mathcal{C} = \{m_1, \dots, m_k\}$ identified via semantic similarity (cosine similarity $> \theta_{cluster}$), the consolidation process performs \textbf{Contrastive Abstraction}:
\begin{enumerate}
    \item \textbf{Cluster Identification}: Group temporally disjoint but semantically coherent memories using DBSCAN on embedding space.
    \item \textbf{Invariant Extraction}: Prompt the LLM to identify recurring patterns while discarding incidental specifics. For example, transforming ``bought coffee on Monday'' and ``ordered latte on Wednesday'' into ``has a regular coffee routine''.
    \item \textbf{Schema Generation}: Create a new semantic memory $m_{semantic}$ with generalized content and inherited metadata.
    \item \textbf{Trace Pruning}: Archive or delete the original episodic traces to reduce redundancy.
\end{enumerate}

Formally, the consolidation operator $\mathcal{T}_{sleep}$ maps a set of episodic memories to a semantic schema:
\begin{equation}
    \mathcal{T}_{sleep}(\mathcal{C}) = \text{LLM}_{abstract}\left(\bigoplus_{m_i \in \mathcal{C}} c_i\right)
\end{equation}
where $\bigoplus$ denotes concatenation with temporal ordering, and $\text{LLM}_{abstract}$ is a prompt-guided abstraction function (see Appendix~\ref{app:prompts}).

\subsection{Meta-Cognitive Adaptation via Bayesian Optimization}
\label{subsec:meta_learning}

A key limitation of existing systems is their use of static, one-size-fits-all parameters. We model memory management as a Markov Decision Process (MDP) where:
\begin{itemize}
    \item \textbf{State} $s_t$: Current memory distribution statistics (mean score, variance, cluster density) and user profile features.
    \item \textbf{Action} $a_t$: Adjustments to cognitive parameters $\{\Delta S_{base}, \Delta \lambda, \Delta \theta_{cluster}\}$.
    \item \textbf{Reward} $r_t$: Implicit user engagement signals, defined as:
    \begin{equation}
        r_t = \alpha_1 \cdot \mathbb{1}[\text{positive feedback}] + \alpha_2 \cdot \text{followup\_depth} - \alpha_3 \cdot \text{correction\_turns}
    \end{equation}
\end{itemize}

Instead of expensive deep RL, we employ \textbf{Gaussian Process-based Bayesian Optimization} (GP-BO) to efficiently explore the parameter space. The meta-learner maintains a posterior distribution over optimal parameters $\phi^*$ for each user:
\begin{equation}
    P(\phi^* | \mathcal{D}_{1:t}) \propto P(r_{1:t} | \phi_{1:t}) \cdot P(\phi^*)
\end{equation}
where $\mathcal{D}_{1:t}$ is the history of parameter-reward pairs. To mitigate cold-start issues, we initialize new users with \textit{population priors} derived from offline clustering of similar user profiles (e.g., ``forgetful'', ``nostalgic'', ``business-efficient'').

The acquisition function uses Expected Improvement (EI) to balance exploration and exploitation:
\begin{equation}
    \phi_{t+1} = \arg\max_\phi \mathbb{E}\left[\max(0, r(\phi) - r(\phi^+))\right]
\end{equation}
where $\phi^+$ is the best parameter configuration observed so far.
